\chapter{指数估计的三区间方法}
\label{app:three-interval}
\appendixsectionlessstatements

本附录的主要目的是证明引理 9.5.29。它是极小曲面型几何偏微分方程问题中经常出现的
一般指数行为的一个原型。\personIgnasiMundetIRiera{}和\personTian{}在 (MT09) 中对伪全纯映射采用了一种
优美的方法来得到这类估计，该方法称为三区间方法。我们部分沿用 (OhZ11a) 在另一
背景下的叙述。同一种论证也可用于其他几何问题中出现的许多情形。这里在长柱面上
伪全纯映射直径很小的当前背景下重述他们的论证。

首先重述引理 9.5.29。

\begin{lemma}
\label{lem:app-b-main-decay}
存在 $L_0>0$ 和 $\epsilon_1>0$，使得对任意 $L\geq L_0$，若
\[
 u\colon[-L-1,L+1]\times S^1\longrightarrow(M,J)
\]
是 $J$-全纯的，且 $\operatorname{diam}(u)\leq\epsilon_1$，则
\[
 \left|\frac{\partial u}{\partial\tau}(\tau,t)\right|
 +\left|\frac{\partial u}{\partial t}(\tau,t)\right|
 \leq Ce^{-\lambda\operatorname{dist}(\tau,\partial[-L-1,L+1])}
\]
对 $\tau\in[-L,L]\times S^1$ 成立，其中 $\lambda>0$，$C>0$ 与
$L\geq L_0$ 无关。事实上，通过令 $\epsilon_1$ 充分小，可以取
$\lambda=2\pi\nu>0$，其中 $\nu=\nu(\epsilon_1)$ 可以任意接近 $1$。
\footnote{译者注：原书把标量 $\tau$ 的取值范围写成
$[-L,L]\times S^1$；按紧前公式，本应由 $(\tau,t)$ 属于该柱面。译文照录。}
\end{lemma}

令 $\iota_g>0$ 是 $M$ 上给定相容度量 $g$ 的\termInjectivityRadius{}。取
$\epsilon>0$，假设对某个 $p\in M$ 有
\[
 u([-L-1,L+1]\times S^1)\subset B_p(\epsilon),
\]
并考虑方程
\[
 \frac{\partial u}{\partial\tau}+J\frac{\partial u}{\partial t}=0.
\]
用指数映射把邻域 $B_p(2\epsilon)$ 等同于 $(T_pM,J_p)$ 中以原点为中心、半径为
$2\epsilon$ 的球后，可以假设 $B_p(2\epsilon)$ 是线性空间 $T_pM$ 的子集。再把
$(T_pM,J_p)$ 等同于带标准度量的 $\mathbb C^n$。$M$ 中的 $J$-全纯映射 $u$ 便化为
$\exp_p^*J$-全纯映射；对 $x\in\mathbb C^n\cong T_pM$，仍将后者记为
$J=J(x)$。

在这一设置下，(MT09) 的观察是：指数估计可化为三个依次相邻的局部 $L^2$ 估计
\begin{equation}
\label{eq:app-b-three-cylinder-local}
 \|du\|_{L^2(Z_{II})}
 \leq\frac12\bigl(\|du\|_{L^2(Z_I)}+\|du\|_{L^2(Z_{III})}\bigr),
\end{equation}
其中 $Z_I,Z_{II},Z_{III}\subset[-L-1,L+1]\times S^1$ 是三个依次相邻的单位长度
柱面，即对某个整数 $i$，它们分别是
\[
 [i-1,i]\times S^1,\qquad[i,i+1]\times S^1,
 \qquad[i+1,i+2]\times S^1.
\]

为了在指数衰减中得到最佳的 $\sigma$（我们需要 $\sigma$ 非常接近 $2\pi$），回忆
(MT09) 的作者定义了常数
\[
 \gamma(c)=\frac1{e^c+e^{-c}}.
\]
$\gamma(c)$ 的重要性在于恒等式
\begin{equation}
\label{eq:app-b-gamma-identity}
 \int_0^1e^{c\tau}\,d\tau
 =\gamma(c)\left[
   \int_{-1}^0e^{c\tau}\,d\tau+
   \int_1^2e^{c\tau}\,d\tau\right],
\end{equation}
它稍后会出现在三个依次相邻的单位长度柱面上 $du$ 的 $L^2$-能量中。注意当
$c>0$ 时，$\gamma(c)$ 是 $c$ 的严格递减函数。

回忆下面的初等引理。

\begin{lemma}[(MT09)，引理 9.4]
\label{lem:app-b-three-interval-sequence}
设 $x_k$（$k=0,1,\ldots,N$）为非负数。若存在固定常数
$\gamma\in(0,1/2)$，使得对所有 $1\leq k\leq N-1$，
\[
 x_k\leq\gamma(x_{k-1}+x_{k+1}),
\]
则
\[
 x_k\leq x_0\xi^{-k}+x_N\xi^{-(N-k)}
\]
对所有 $1\leq k\leq N-1$ 成立，其中
\[
 \xi=\frac{1+\sqrt{1-4\gamma^2}}{2\gamma}.
\]
\end{lemma}

\begin{remark}
若 $\gamma=\gamma(c)=(e^c+e^{-c})^{-1}$，则可验证 $\xi=e^c$，上面的不等式成为
指数衰减估计
\begin{equation}
\label{eq:app-b-sequence-decay}
 x_k\leq x_0e^{-ck}+x_Ne^{-c(N-k)},
\end{equation}
其中 $1\leq k\leq N-1$。
\end{remark}

\begin{proposition}
\label{prop:app-b-local-four-pi}
对任意 $0<\nu<1$，存在 $L_0>0$ 和只依赖于 $\nu$、但与 $L\geq L_0$ 无关的
$\epsilon_1=\epsilon_1(\nu)>0$，使得在一切满足
\[
 u\colon[-L-1,L+1]\times S^1\longrightarrow(M,J)
\]
为 $J$-全纯、$\operatorname{diam}(u)\leq\epsilon$、
$0<\epsilon<\epsilon_1$ 且 $p\in M$ 的情形下，对
$[-L,L]\times S^1$ 中任意三个依次相邻的柱面 $Z_I,Z_{II},Z_{III}$，都有
\begin{equation}
\label{eq:app-b-gamma-four-pi}
 \|du\|_{L^2(Z_{II})}
 \leq\gamma(4\pi\nu)
 \bigl(\|du\|_{L^2(Z_I)}+\|du\|_{L^2(Z_{III})}\bigr).
\end{equation}
\end{proposition}

\begin{proof}
用反证法证明\eqref{eq:app-b-gamma-four-pi}。假设存在序列 $L_j\to\infty$ 和
$u_j$，满足
\[
 \operatorname{diam}\bigl(u_j([-L_j,L_j]\times S^1)\bigr)\longrightarrow0,
\]
但在 $[-L_j,L_j]\times S^1$ 中三个依次相邻的柱面
$Z_I^j,Z_{II}^j,Z_{III}^j$ 上违反\eqref{eq:app-b-gamma-four-pi}：
\[
 \|du_j\|_{L^2(Z_{II}^j)}
 >\gamma(4\pi\nu)
 \bigl(\|du_j\|_{L^2(Z_I^j)}+\|du_j\|_{L^2(Z_{III}^j)}\bigr).
\]
在 $Z_I^j\cup Z_{II}^j\cup Z_{III}^j$ 上考察重标度序列
\[
 \widehat u_j=
 \frac{u_j}{\|u_j\|_{L^\infty(Z_I^j\cup Z_{II}^j\cup Z_{III}^j)}}.
\]
分母绝不为 $0$；否则 $u_j\equiv0$ 于这三个柱面的并上，与 $u_j$ 非常值的假设矛盾。
于是
\[
 \|\widehat u_j\|_{L^\infty(Z_I^j\cup Z_{II}^j\cup Z_{III}^j)}=1,
\]
\[
 \|d\widehat u_j\|_{L^2(Z_{II}^j)}
 >\gamma(4\pi\nu)
 \bigl(\|d\widehat u_j\|_{L^2(Z_I^j)}
       +\|d\widehat u_j\|_{L^2(Z_{III}^j)}\bigr),
\]
且
\[
 \left(\frac\partial{\partial\tau}
 +J(u_j)\frac\partial{\partial t}\right)\widehat u_j=0.
\]

适当平移 $Z_I^j\cup Z_{II}^j\cup Z_{III}^j$，并回忆
$\operatorname{diam}(u_j)\to0$，可选取一个仍记作 $u_j$ 的子列，使得在
$[0,3]\times S^1$ 上：
\begin{enumerate}[label=(\arabic*)]
\item 重标度映射 $\widehat u_j(\tau,t)$ 收敛到 $\widehat u_\infty(\tau,t)$；
\item 原映射 $u_j(\tau,t)$ 收敛到某点 $p_\infty\in M$。
\end{enumerate}
下面还需要在 $j\to\infty$ 时 $J(u_j(\tau,t))$ 的 $C^1$ 行为收敛到 $C^1$ 行为。

\begin{lemma}
\label{lem:app-b-derivative-vanishing}
必要时取子列，则对任意给定紧子集 $K\subset\mathbb R\times S^1$，有
\[
 \lim_{j\to\infty}\|du\|_{C^0;K}=0.
\]
特别地，$J(u_j)\to J_0$ 于 $\mathbb R\times S^1$ 上的紧 $C^1$ 拓扑。
\footnote{译者注：极限左端原书写作 $\lVert du\rVert_{C^0;K}$，漏去序列下标
$j$；译文照录原式。}
\end{lemma}

\begin{proof}
给定紧子集 $K\subset\mathbb R\times S^1$。由于 $L_j\to\infty$，最终有
$K\subset[-L_j,L_j]\times S^1$。反设必要时取子列后，存在某个 $\epsilon>0$ 和
$(\tau_j,t_j)\in K$，使得对所有 $j$，
$|du_j(\tau_j,t_j)|>\epsilon$。考虑平移并旋转后的映射
\[
 \widetilde u_j(\tau,t)=u_j(\tau-\tau_j,t-t_j),
\]
其定义域包含 $[-L_j+\tau_j,L_j+\tau_j]\times S^1$，且
$|d\widetilde u_j(0,0)|>\epsilon$。标准冒泡论证要么产生一个非常值气泡，要么产生
一个子列，使 $\widetilde u_j\to u_\infty$ 于紧 $C^1$ 拓扑，其中
$u_\infty\colon\mathbb R\times S^1\to M$ 是满足
\[
 |du_\infty(0,0)|\geq\epsilon
\]
的 $J_0$-全纯映射，因而不是常值映射。这导致矛盾，因为
$\operatorname{im}\widetilde u_j=\operatorname{im}u_j\subset B_{p_j}(2\epsilon_j)$，
所以 $\operatorname{diam}(\operatorname{im}u_\infty)=0$。引理得证。
\end{proof}

特别地，该引理蕴含 $J(u_j(\tau,t))\to J_0$ 于\emph{固定}柱面
$[0,3]\times S^1(=Z_I\cup Z_{II}\cup Z_{III})$ 上的 $C^1$ 拓扑。因此由以上讨论，
极限映射 $\widehat u_\infty$ 满足
\[
 \|\widehat u_\infty\|_{L^\infty(Z_I\cup Z_{II}\cup Z_{III})}=1,
\]
\[
 \|d\widehat u_\infty\|_{L^2(Z_{II})}
 \geq\gamma(4\pi\nu)
 \bigl(\|d\widehat u_\infty\|_{L^2(Z_I)}
       +\|d\widehat u_\infty\|_{L^2(Z_{III})}\bigr),
\]
以及
\[
 \left(\frac\partial{\partial\tau}
 +J_0\frac\partial{\partial t}\right)\widehat u_\infty=0.
\]
由第一、第三个恒等式，$\widehat u_\infty$ 是到 $\mathbb C^n$ 的非零全纯映射。把它
展开为傅里叶级数
\[
 \widehat u_\infty(\tau,t)
 =\sum_{k=-\infty}^{\infty}a_ke^{2\pi k\tau}e^{2\pi kit},
 \qquad
 \|\widehat u_\infty\|_{L^2(Z_I\cup Z_{II}\cup Z_{III})}\leq3,
\]
其中 $a_k$ 是 $\mathbb C^n$ 中的常向量。直接计算得
\begin{equation}
\label{eq:app-b-fourier-energy}
 \|d\widehat u_\infty\|_{L^2([a,b]\times S^1)}^2
 =\sum_{k=-\infty}^{\infty}4\pi^2k^2|a_k|^2
   \int_a^be^{4\pi\tau}\,d\tau.
\end{equation}
把\eqref{eq:app-b-gamma-identity}乘以 $e^{4\pi k}$，并在其中令 $c=4\pi$，得到
\begin{equation}
\label{eq:app-b-shifted-gamma-identity}
 \int_k^{k+1}e^{4\pi\tau}\,d\tau
 =\gamma(4\pi)\left[
   \int_{k-1}^ke^{4\pi\tau}\,d\tau
   +\int_{k+1}^{k+2}e^{4\pi\tau}\,d\tau\right].
\end{equation}
由\eqref{eq:app-b-fourier-energy}与\eqref{eq:app-b-shifted-gamma-identity}可见
\[
 \|d\widehat u_\infty\|_{L^2(Z_{II})}
 =\gamma(4\pi)
 \bigl(\|d\widehat u_\infty\|_{L^2(Z_I)}
       +\|d\widehat u_\infty\|_{L^2(Z_{III})}\bigr).
\]
这与
\[
 \|d\widehat u_\infty\|_{L^2(Z_{II})}
 \geq\gamma(4\pi\nu)
 \bigl(\|d\widehat u_\infty\|_{L^2(Z_I)}
       +\|d\widehat u_\infty\|_{L^2(Z_{III})}\bigr)
\]
矛盾，因为 $\gamma(4\pi)<\gamma(4\pi\nu)$。
\footnote{译者注：按紧前傅里叶展开，模态 $k$ 的能量指数通常应含 $k$；原书
\eqref{eq:app-b-fourier-energy} 却印作 $e^{4\pi\tau}$，继而用与傅里叶模态无关的
$\gamma(4\pi)$ 得到等式。译文照录这一推导，不自行重构。}
\end{proof}

结合上面的引理与\eqref{eq:app-b-sequence-decay}，得到以下推论。

\begin{corollary}
\label{cor:app-b-energy-decay}
对任意 $0<\nu<1$，存在 $L_0>0$ 和只依赖于 $\nu$、但与 $L\geq L_0$ 无关的
$\epsilon_0=\epsilon_0(\nu)>0$，使得对所有 $0<\epsilon\leq\epsilon_1$ 及
$[\tau,\tau+1]\subset[-L,L]$，都有
\[
 \int_{[\tau,\tau+1]\times S^1}|du|^2
 \leq e^{-4\pi\nu(L-|\tau|)}
 \left[
  \int_{[-L-1,-L]\times S^1}|du|^2
  +\int_{[L,L+1]\times S^1}|du|^2
 \right].
\]
\footnote{译者注：本推论先引入 $\epsilon_0=\epsilon_0(\nu)$，量词中却仍使用
$\epsilon_1$；译文照录这两个下标。}
\end{corollary}

由这些结果以及柱面 $[\tau-\tfrac12,\tau+\tfrac12]\times S^1$ 上的标准椭圆估计，
得到下面关于 $u$ 的逐点指数衰减估计。

\begin{corollary}
\label{cor:app-b-pointwise-decay}
对任意 $0<\nu<1$，存在 $L_0>0$ 和只依赖于 $\nu$、但与 $L\geq L_0$ 无关的
$\epsilon_1=\epsilon_0(\nu)>0$，使得对所有 $0<\epsilon\leq\epsilon_1$ 及
$\tau\in[-L,L]$，都有
\[
 |du|\leq Ce^{-2\pi\nu(L-|\tau|)}
 \bigl(
  \|du\|_{L^2([-L-1,-L]\times S^1)}
  +\|du\|_{L^2([L,L+1]\times S^1)}
 \bigr).
\]
$\nu$ 可以任意接近 $1$。
\end{corollary}

\begin{remark}
同一方案可应用于半带 $[0,\infty)\times[0,1]$ 上的弗洛尔方程
\[
 \begin{cases}
 \partial u/\partial\tau+J\,\partial u/\partial t=0,\\
 u(\tau,0)\in L_0,\quad u(\tau,1)\in L_1,
 \end{cases}
\]
其渐近极限为
$\lim_{\tau\to\infty}u(\tau,\cdot)=p\in L_0\cap L_1$，这里 $p$ 是横截交点；此时
极限方程为
\[
 \begin{cases}
 \partial\widehat u_\infty/\partial\tau
 +J(p)\,\partial\widehat u_\infty/\partial t=0,\\
 \widehat u_\infty(\tau,0)\in T_pL_0,\quad
 \widehat u_\infty(\tau,1)\in T_pL_1.
 \end{cases}
\]
这给出了第 14.1 节所证指数衰减结果的另一种证明；它不仅计算较少，还给出了
式 (8.6.51) 的指数中 $\delta$ 的精确值。
\end{remark}

\begin{exercise}
补全上述证明的细节。
\end{exercise}
